Aquí está el plan detallado, comentario por comentario, con las acciones concretas que haría en el artículo:

Comentario 1 — Motivación bien articulada, pero limitaciones importantes

Acción: No requiere cambios estructurales, pero sí reforzar el abstract y la introducción para anticipar proactivamente las limitaciones y enmarcarlas como trabajo en progreso justificado. Añadiría una frase explícita tipo:

"While HAMELIN does not introduce new modeling algorithms, its contribution lies in the protocol-aware interface design and the operationalization of HGML principles in a clinical context — an area where no prior system exists."

Comentario 2 — Contribución es de diseño, no algorítmica

Este es el comentario más delicado. El revisor tiene razón, pero la contribución sí es válida si se enmarca correctamente. Haría lo siguiente:

2a. Reformular la contribución en la introducción con una lista explícita de contribuciones (contributions paragraph), algo habitual en papers IEEE:

"The contributions of this paper are: (1) a protocol-aware interface design grounded in real clinical protocols; (2) a mapping between clinical research concepts and ML pipeline stages; (3) an operationalization of HGML criteria in a clinical AutoML system; (4) a participatory design methodology involving clinical research professionals."

2b. Añadir una comparativa con herramientas existentes en la sección de Related Work, en forma de tabla, mostrando que ninguna herramienta existente cubre las dimensiones que HAMELIN cubre (protocol-awareness, HGML alignment, clinical terminology, inclusion/exclusion criteria support). Esto refuerza la novedad sin necesidad de contribución algorítmica.

Feature	AutoGluon	H2O	Azure AutoML	HAMELIN
Protocol-aware	✗	✗	✗	✓
Inclusion/exclusion criteria	✗	✗	✗	✓
Clinical terminology	✗	✗	Partial	✓
HGML alignment	✗	✗	✗	Partial
No-code for clinicians	Partial	Partial	✓	✓

2c. En la sección de Discussion, añadir un párrafo que defienda explícitamente el valor de las contribuciones de diseño en sistemas HCI/HCAI, citando trabajos que avalen este tipo de contribución (e.g., trabajos de Shneiderman sobre HCAI, o papers de CHI/CSCW que son puramente de diseño).

Comentario 3 — No hay estudio de usuario formal

Este es el comentario más difícil de resolver si no se tiene tiempo de hacer un estudio nuevo. Hay dos opciones:

Opción A (ideal): Realizar un estudio de usabilidad estructurado con los CRPs ya involucrados, usando un cuestionario estándar como el SUS (System Usability Scale) o el NASA-TLX. Con 4-5 participantes ya es publicable en un paper de diseño de sistemas. Presentar los resultados en una nueva sección o subsección.

Opción B (si no hay tiempo): Reconvertir la evidencia existente en una evaluación experta estructurada. Documentar las sesiones de feedback con los CRPs de forma sistemática: número de sesiones, participantes, tipo de feedback recibido, cambios realizados. Añadir una tabla o figura que muestre las iteraciones de diseño y cómo el feedback se tradujo en cambios concretos. Esto no es un estudio de usuario formal, pero es evidencia sistemática y honesta.

En cualquier caso, añadir en Limitations una frase que reconozca explícitamente la ausencia de un estudio de usuario formal y lo proponga como trabajo futuro.

Comentario 4 — Diseño participativo anecdótico, no sistemático

Complementario al comentario 3. Acciones concretas:

4a. Añadir una tabla de iteraciones de diseño en la sección 3.1 (Participatory Design Process) con columnas como: Iteration, Protocol used, Participants, Key findings, Changes made. Esto convierte la narrativa anecdótica en evidencia estructurada.

4b. Describir el proceso con más rigor metodológico: número exacto de sesiones, duración, perfil de los participantes (sin datos personales), método de recogida de feedback (notas, grabaciones, formularios), y criterios de convergencia del diseño.

4c. Si se usaron los mockups para validación, describir explícitamente qué se validó, cómo se recogió el feedback y qué cambios resultaron. La Figura 2 (mockup) ya está, pero necesita más contexto narrativo.

Comentario 5 — No sigue las guías de formato IEEE

Acciones de formato a revisar sistemáticamente:

5a. Estructura del documento: Verificar que el orden de secciones sigue el estándar IEEE: Abstract → Keywords → Introduction → Background/Related Work → Methodology → Results/Evaluation → Discussion → Conclusions → References.

5b. Figuras y tablas:

Las figuras deben tener numeración y captions en la parte inferior.
Las tablas deben tener captions en la parte superior (actualmente la Tabla 1 parece correcta, pero verificar).
Las figuras [h], [H], [htbp] deben revisarse; IEEE prefiere [t] o [b] para posicionamiento.

5c. Referencias: IEEE usa numeración entre corchetes [1], [2]... El estilo unsrt de BibTeX es compatible, pero verificar que no haya referencias con formato incorrecto (URLs sin formato, años faltantes, etc.).

5d. Columnas y márgenes: Verificar que el documento usa "documentclass[conference]{IEEEtran}" correctamente y que no hay elementos que rompan el layout de doble columna (las figuras figure* son correctas para ancho completo, pero verificar que no causen problemas de posicionamiento).

5e. Paquetes no estándar: tikz, booktabs, multirow, float son aceptables, pero \scalebox dentro de figuras puede causar problemas de escala en la versión final. Considerar rediseñar las figuras TikZ con tamaños absolutos en lugar de escalado relativo.

5f. Sección de agradecimientos: IEEE requiere una sección \section*{Acknowledgment} antes de las referencias, no un \thanks{} en el título (aunque \thanks también es válido en IEEEtran, es preferible la sección explícita para visibilidad).

5g. Comentarios % Revisar: Eliminar todos los comentarios internos del código LaTeX antes de la sumisión final.

Resumen de prioridades
Prioridad	Comentario	Esfuerzo estimado
�� Alta	C3: Añadir evaluación estructurada con CRPs	Alto
�� Alta	C2: Añadir tabla comparativa y párrafo de contribuciones	Medio
�� Media	C4: Estructurar el proceso participativo con tabla de iteraciones	Medio
�� Media	C5: Revisión de formato IEEE	Bajo-Medio
�� Baja	C1: Reforzar framing en abstract/intro	Bajo

¿Quieres que empiece a redactar los cambios concretos en el LaTeX para alguno de estos puntos?